\documentclass[11pt]{article}
\usepackage[utf8]{inputenc}
\usepackage[T1]{fontenc}
\usepackage{lmodern}
\usepackage[margin=1in]{geometry}
\usepackage{amsmath,amssymb}
\usepackage{textcomp}
\usepackage[dvipsnames]{xcolor}
\usepackage{enumitem}
\usepackage{parskip}
\usepackage{titlesec}
\usepackage[colorlinks=true,linkcolor=blue,citecolor=blue,urlcolor=blue]{hyperref}
\usepackage[round]{natbib}
\usepackage{microtype}
\emergencystretch=1.5em

\titleformat{\section}{\large\bfseries}{\thesection}{1em}{}
\titleformat{\subsection}{\normalsize\bfseries}{\thesubsection}{1em}{}
\titleformat{\subsubsection}{\normalsize\bfseries\itshape}{}{0em}{}

\title{\textbf{Recommending Ghosts: Install-Time Provenance Gating of\\ Squatted Skill Recommendations in Live Agent Trajectories}}
\author{}
\date{}

\newcommand{\ph}{\textit{[placeholder --- fill in after experiments are run]}}

\begin{document}
\maketitle
\vspace{-2em}

\begin{quote}
\small\textit{Assembled from the six confirmed stage outputs of an \texttt{auto-research-discussion} run (seed topic B1.6, autoclaude self-play), post research-taste audit (\textcolor{ForestGreen}{\textbf{PASS}}) and stage review (three REVISE findings applied). Status: \textbf{specification / paper draft --- no experiments have been run.} All numbers below are pre-registered design targets with stated rationale, not results.}
\end{quote}

\begin{abstract}
Autonomous LLM agents increasingly recommend and auto-install skills and tools from open registries, and a just-published large-scale study shows they hallucinate skill names that do not exist at rates of 36--43\% across every tested configuration --- a confirmed supply-chain attack precondition, because an adversary who observes which fake names agents reliably produce can pre-register a malicious skill under one of them and wait for a victim's agent to recommend and install it. Every mitigation in that study acts at the recommendation step, by constraining what the model is allowed to say, and every one buys a lower hallucination rate only by destroying usefulness (best case: a real-skill recall of roughly 16\%). We observe that all of these defenses share one structural property --- they intervene on the recommendation distribution --- and that a defense placed one step later, at install time, need not touch that distribution at all: it need only ask whether the specific skill about to be installed comes from a trusted publisher. Crucially, a bare \emph{existence} check cannot substitute, because the attack is squatting: once the malicious skill is registered, its name exists. We build a multi-step agent-trajectory testbed that embeds a recommend-then-install decision inside a task with a real downstream outcome, insert a gate keyed only on publisher-trust-tier provenance metadata (no name lookup, no name-similarity heuristic, no model judgment), seed a squatting adversary from the confirmed recurring-name pattern, and measure the resulting (compromise-block, legitimate-reject) operating point, how it moves as interception is delayed, and the recovery cost --- reporting all of it per model and harness against the published recommendation-time defenses as fixed reference markers, not competitors to beat.
\end{abstract}

\section{Introduction}

Autonomous LLM agents no longer just call a fixed set of built-in tools; they discover, recommend, and auto-install skills from open registries mid-task, on behalf of a principal who never sees the individual install decision. A just-published large-scale security study settles an uncomfortable fact about this practice: across 15,000 prompts and twelve model-and-harness configurations, every single configuration recommends skills that do not exist, at an average rate of 36--43\%, and these hallucinated names recur predictably enough across prompts and models to be reliable attack targets \citep{hallucskills}. That recurrence is what turns a reliability quirk into a supply-chain vulnerability: an adversary who watches which fake names an agent reliably produces can pre-register a real, malicious skill under one of those names on a registry that rarely verifies publishers, then wait for a victim's agent to recommend and auto-install the trap. The study confirms the full path end to end in a sandboxed proof of concept. The attack-precondition pattern itself --- hallucinated dependency names becoming squattable attack surface --- was established earlier for package ecosystems, where code-generating LLMs hallucinate package names at scale and ``slopsquatting'' attacks pre-register them \citep{spracklen,phantoms}, and where a 2026 replication finds the threat persisting across frontier models, with dozens of model-agnostic hallucinated names still registrable after registries' deployed defenses \citep{packrange}. What is different in the skill setting --- and what this paper's question turns on --- is the consumer of the hallucinated name: not a developer reading a suggestion, but an agent that autonomously installs and executes mid-task.

A just-published study settling the phenomenon leaves one intervention point unexamined. Every mitigation that study evaluates --- retrieval-constrained generation, self-refinement, cross-model voting, and their combination --- acts at the \emph{recommendation} step, constraining what the model is allowed to say, and every one of them buys a lower hallucination rate only by collapsing usefulness: the strongest cuts hallucination from 40.8\% to 3.2\% but keeps only about 16\% of the correct real skills \citep{hallucskills}. This is not four independent failures; it is one structural property. A filter operating on the model's output cannot tell a real but training-atypical skill name apart from a hallucinated one, because on the recommendation distribution they look the same. The study's own closing sentence calls for ``registry-level name reservations and verified recommendation pipelines,'' and builds neither.

We observe that the checkpoint does not have to sit on the recommendation distribution at all. A defense placed one step later, at the moment of install, never needs to constrain what the model says; it needs only to ask whether the specific skill about to be installed comes from a trusted publisher. Moving the checkpoint this way separates two signals the recommendation-time filters necessarily conflate: a correctly-named real skill from an unverified publisher, and a hallucinated name, are different things to an install-time provenance check even though they are indistinguishable to a recommendation-time one. And crucially, a \emph{bare existence} check --- does a skill by this name exist in the registry --- cannot do this job, because the attack is squatting: the whole point is that the attacker registers the hallucinated name, so by the time a victim's agent recommends it, the name exists and a naive existence gate waves the malicious skill through. The signal that survives the threat model is publisher provenance, not name existence. We do not claim the mechanism class as ours: registry-side publisher-identity verification and trust-scored tool registries have been proposed, at design level, for exactly the tool-squatting threat \citep{ztr,etdi}, and package ecosystems deploy provenance attestation in practice. What no prior work provides is the measurement --- whether such a gate actually blocks squatting seeded from confirmed hallucinated-name recurrence, what it costs on legitimate skills from unverified publishers, and whether interception timing matters. That measurement, not the gate, is this paper's contribution.

Turning this into a measurable question surfaces its own difficulties. The evaluation's ground truth must make squatting real: if ``the name exists'' is treated as safe, the trap is scored safe by construction, so the registry snapshot must carry forgery-resistant publisher-trust tiers and the simulated adversary must pre-register names drawn from the \emph{confirmed recurring} hallucinated set, not arbitrary strings. The gate itself must reach its verdict without any model judgment, because the same study shows self- and cross-model auditing separates real from hallucinated names no better than a coin flip (51--53\%); an LLM ``looks-real'' call inside the gate would import exactly that failure. The usability cost must be measured on \emph{real} skills --- a legitimate first-time publisher is indistinguishable from a squatter on provenance alone until trust accrues, so whether the gate avoids or merely relocates the recall collapse is an empirical question, not an assumption. And \emph{when} the check fires is itself a first-class outcome, not a tuning detail: unlike a hallucinated non-existent name, which stalls the task and surfaces organically because the tool is simply not there, a squatted skill that installs and runs can complete the visible task while doing its damage silently, so a single post-hoc audit of the final state cannot tell ``blocked before execution'' apart from ``executed silently, task succeeded anyway.'' Only step-indexed interception distinguishes the two, and the number of steps between an untrusted install and its detection is the recovery cost an operator pays.

Our approach turns each of these obstacles into a specific design choice rather than a hedge. The registry snapshot is annotated with publisher-trust tiers and registration recency, and the squatting adversary is seeded from the confirmed recurring-name pattern, so squatting is the ground truth before any agent runs. The gate keys on provenance metadata alone --- no existence lookup as a safety signal, no name-similarity heuristic, no model judgment --- so its block decisions cannot inherit either the name-shortcut brittleness or the coin-flip detectability the base study documents; the recurring-name set is consumed only by the offline adversary and the ground-truth labeler, never read by the online gate. The usability cost is measured directly, as the rejection rate on a held-out probe set of real, useful skills from unverified publishers. And interception timing is manipulated by an explicit delayed-gate control that holds the trajectory and task difficulty fixed while varying only how many steps late the check fires, so timing's effect on the security outcome is isolated from task difficulty.

We evaluate this design across multi-step trajectories that embed a recommend-then-install decision, drawn from the base study's own two settings --- standalone-LLM recommendation and agent-harness recommend-then-install --- comparing a gate-on arm against gate-off, a bare-existence placebo gate, and a delayed-gate control, and reporting every operating point and timing statistic per model and harness against the base study's four published recommendation-time defenses as fixed reference markers on the same plane.

In summary, we contribute:
\begin{itemize}[leftmargin=1.4em]
\item An identification of the structural property shared by all recommendation-time defenses against skill hallucination --- they constrain the recommendation distribution, which is why they collapse recall --- and the analysis showing why, under the base study's own findings, an install-time checkpoint must key on publisher provenance rather than name existence, since the squatting attack renders existence-checking inert by construction. The checkpoint mechanism itself is conceded to design-level registry-trust proposals \citep{ztr,etdi}; the causal analysis, and the measurement of what the mechanism class actually buys and costs, are not provided by any of them.
\item A multi-step recommend-then-install trajectory testbed with a trust-tiered registry snapshot and a squatting adversary seeded from confirmed hallucinated-name recurrence, faithful to the base study's measured prevalence and targetability, released as the instrument for this and future install-time interventions.
\item A pre-registered measurement of the (compromise-block, legitimate-unverified-publisher-reject) operating point that install-time provenance gating achieves, placed on the same plane as the published recommendation-time defenses, together with a detection-timing-to-cost curve and a falsification control (a bare-existence placebo gate) that isolates whether provenance, not mere existence, is doing the work.
\end{itemize}

\section{Problem Definition}

\noindent\textbf{Formal setup.} A task trajectory is a sequence of steps $\tau = (x_1, \dots, x_L)$ containing one recommend-then-install step $x_r$ that proposes a skill name $n$ and, if the install proceeds, fetches and executes a skill body $b$ drawn from a registry $R$. A registry snapshot $R = \{(n_k, p_k, \theta_k)\}$ pairs each published name $n_k$ with its publisher $p_k$ and a publisher-trust tier $\theta_k \in \{\mathrm{verified}, \mathrm{known}, \mathrm{unknown}\}$ together with registration recency. A gate $g_\delta$ maps the proposed name and the snapshot to a verdict, $g_\delta : (n, R) \to \{\mathrm{allow}, \mathrm{block}, \mathrm{review}\}$, and fires $\delta$ steps after $x_r$: $\delta = 0$ is immediate interception, $\delta = k$ is a delayed check, and $\delta = \infty$ is no gate at all (gate-off). A squatting adversary pre-registers a subset of the confirmed recurring hallucinated names $\mathcal{H}$ under an unverified publisher identity ($\theta = \mathrm{unknown}$); a disjoint probe set $\mathcal{L}$ holds real, useful skills from legitimate but unverified publishers. The object of study is the map, per model and harness cell, from an intervention $(g, \delta)$ to four quantities: the rate at which a squatted install on an $\mathcal{H}$-name is completed and executes (from sandbox state), the rate at which a legitimate skill in $\mathcal{L}$ is rejected, the downstream task success, and the recovery cost --- the number of extra steps and tokens between an untrusted install and task resolution.

\noindent\textbf{Why this is hard.} Four difficulties are intrinsic to the setting. First, making the measurement faithful to the attack requires that ``the name exists'' be true for the malicious skill: because the adversary registers an $\mathcal{H}$-name, a ground truth keyed on name existence would score the trap safe by construction, so the snapshot must carry a distinct provenance predicate $\theta$ that existence cannot satisfy. Second, the verdict $g$ must be a function of $\theta$ and recency alone --- no model judgment --- because self- and cross-model auditing of $n$ is known to separate real from hallucinated names no better than chance, so any LLM inside the gate re-imports that failure. Third, the usability cost is the rejection rate on $\mathcal{L}$, and a legitimate first-time publisher occupies exactly the same $\theta = \mathrm{unknown}$ tier as a squatter until trust accrues, so whether the gate avoids or merely relocates the recall collapse cannot be assumed --- it must be measured against real skills. Fourth, the outcome on a squatted install is security-relevant and, unlike a stalled task, need not surface organically: a skill that installs and executes can leave the visible task successful, so distinguishing ``blocked before execution'' from ``executed silently'' requires manipulating $\delta$ under a fixed trajectory rather than reading the final state.

\section{Method Overview}

At a conceptual level the study is a single measurement pipeline with six cooperating parts, three of which construct a faithful testbed offline and three of which run the intervention and account for its effect online. A trust-tiered registry snapshot and a recurrence-seeded squatting adversary make squatting the ground truth before any agent runs; a recommend-then-install trajectory harness embeds the install decision inside a real downstream task; a provenance-only install gate reaches a verdict from publisher-trust metadata alone; a detection-timing instrument records and manipulates how many steps late the check fires; and an operating-point accounting turns every episode into a point on the shared (compromise-block, legitimate-reject) plane, reported per model and harness.

The \textbf{trust-tiered registry snapshot} annotates each published skill with a publisher-trust tier and registration recency, so a squatted malicious skill can satisfy ``exists'' while failing ``verified-or-known publisher.'' The \textbf{recurrence-seeded squatting adversary} pre-registers names drawn from the confirmed recurring hallucinated-name set rather than arbitrary strings, keeping the attack faithful to the demonstrated one. The \textbf{recommend-then-install trajectory harness} embeds one install decision inside a longer task whose outcome carries a security dimension --- did a squatted skill install and execute --- alongside task success. The \textbf{provenance-only install gate} reaches its verdict from the snapshot's trust tier alone, never reading the recommended name's string, so it inherits neither the name-shortcut brittleness nor the coin-flip detectability of the recommendation-time defenses. The \textbf{detection-timing instrument} records the step index at which an untrusted install is first caught and provides a delayed-gate control that varies only that timing. The \textbf{operating-point accounting} records, per episode, whether a compromise was blocked or completed, whether the task succeeded, the recovery cost, and the rejection rate on the legitimate probe set, reporting a single operating point per cell against the base study's four published points.

\section{Methodology}

\subsection{Trust-Tiered Registry Snapshot}
\textbf{Challenge.} Because the adversary registers a name the agent reliably hallucinates, ``a skill by this name exists'' is true for the malicious skill; an evaluation whose ground truth is name existence would score the trap safe by construction.

\textbf{Insight.} Existence and provenance are different predicates, and only the second survives the squatting threat model --- so the testbed must carry provenance, not merely a name list.

\textbf{Method.} Each published name in the snapshot is annotated with a publisher-trust tier --- verified (a publisher that passed an identity/verification step), known (a publisher previously seen and trusted by this agent or principal), or unknown (anonymous or first-seen) --- plus registration recency. This is an offline property of the testbed, fixed before any agent runs: a squatted malicious skill occupies the unknown tier, so it satisfies existence while failing the verified-or-known predicate. The snapshot is the pre-committed ground truth for whether a given install is a compromise, independent of any model judgment.

\subsection{Provenance-Only Install Gate}
\textbf{Challenge.} A verdict that consulted a model's ``does this name look real'' judgment would inherit the confirmed coin-flip (51--53\%) detectability of hallucinated names, and a verdict that keyed on name existence or name similarity would inherit the documented tendency of models to shortcut on names rather than semantics.

\textbf{Insight.} Moving the checkpoint to the install step means the verdict never has to touch the recommendation distribution --- it can be a pure function of who published the skill --- so it does not pay the recall cost the recommendation-time filters pay, and it can exclude both the model judgment and the name signal that the base study shows are unreliable.

\textbf{Method.} The gate reads only the snapshot's trust tier for the proposed skill's publisher, never the recommended name's string: allow iff the publisher is verified or known; otherwise block (unknown publisher) or route to human-review (an operator-configurable middle verdict for an unknown publisher an operator may wish to approve manually). No existence lookup is used as a safety signal, no similarity-to-known-hallucinations heuristic, and no model judgment enters the verdict path; the confirmed recurring-name set is consumed only by the offline adversary and the ground-truth labeler, never by this online gate. The human-review verdict is pinned onto both reporting axes before any downstream measurement uses it: on a legitimate probe skill it counts as a rejection for the conservative (worst-case usability) reading of the reject-rate axis, and is \emph{also} reported as a separate deferred-not-rejected count, so the usability cost is given as a bounded interval --- block-only at the optimistic end, block-plus-review at the conservative end --- rather than a single number that hides the human-review policy; on a squatted skill it counts as blocked, since the security claim is about autonomous auto-install, and a skill held for a human is not autonomously installed.

\subsection{Recommend-Then-Install Trajectory Harness}
\textbf{Challenge.} A single-shot recommendation event cannot expose the difference between an install blocked before execution and one that executes silently while the visible task still succeeds --- the distinction the security outcome turns on.

\textbf{Insight.} Only an install decision embedded inside a longer task with a real, checkable downstream outcome makes that distinction observable, because only then does ``the task succeeded'' fail to imply ``nothing malicious ran.''

\textbf{Method.} One recommend-then-install decision is embedded as a single step inside a longer, executable multi-step task with a checkable outcome, instantiated from the base study's two settings --- standalone-LLM recommendation and agent-harness recommend-then-install. Each episode records a two-part outcome: task success (did the trajectory complete its goal) and a security label (did a squatted skill install and execute, read from sandbox state), so the two can diverge.

\subsection{Recurrence-Seeded Squatting Adversary}
\textbf{Challenge.} An adversary that pre-registered arbitrary malicious names would over- or under-state the gate's value, because the demonstrated attack depends on the specific property that hallucinated names recur predictably enough to be cheaply squatted.

\textbf{Insight.} Faithfulness to the measured attack requires seeding the adversary from the confirmed recurrence, not from a hypothetical worst case --- the recurrence is precisely what makes pre-registration low-effort and the threat real.

\textbf{Method.} The simulated attacker pre-registers, under the unknown-publisher tier, a subset of the confirmed recurring hallucinated names elicited by the base study's own prompt-construction method, rather than arbitrary strings. This keeps the attack path identical to the demonstrated one and lets the block-rate axis measure blocking of genuinely squatted names. The recurring-name set is used here (offline, to build the attack and label ground truth) and nowhere else --- in particular the online gate never sees it.

\subsection{Detection-Timing Instrumentation}
\textbf{Challenge.} The causal effect of interception timing on a security outcome cannot be read from observed trajectories, because the timing at which an install is caught is confounded with how hard the task was and with whether the trajectory failed for unrelated reasons.

\textbf{Insight.} A benign convergence of ``caught early'' with ``easy task'' can be broken by an intervention that holds the trajectory and task difficulty fixed and varies only the step at which the check fires --- a controlled delay, not an observed one.

\textbf{Method.} The instrument records the step index at which an untrusted install is first caught: immediately by the gate ($\delta = 0$), never or late via organic task/security failure (gate-off, $\delta = \infty$), or artificially delayed by $k$ steps (the delayed-gate control). The delayed-gate arm holds the trajectory and task difficulty fixed and moves only $\delta$, isolating timing's effect on the security outcome and on recovery cost. This controlled-delay design is original; it borrows the counterfactual logic of holding everything fixed but the intervention point from prior failure-attribution work \citep{agentracer}, not that work's offline post-completion tracer.

\subsection{Operating-Point Accounting}
\textbf{Challenge.} Whether the gate solved the security/usability tradeoff or merely relocated it cannot be read off a block rate alone, and a pooled number would be uninterpretable against the base study's own per-configuration reporting.

\textbf{Insight.} The single number that decides ``solved versus relocated'' is the rejection rate on \emph{real} skills; reported alongside the block rate as one point on a shared plane, per model and harness, it becomes directly comparable to the published recommendation-time defenses.

\textbf{Method.} Per episode, the accounting records whether a compromise was blocked or completed, whether the task succeeded, the recovery cost (extra steps and tokens between an untrusted install and task resolution), and the rejection rate on the legitimate-unverified-publisher probe set. Each model-and-harness cell yields a single (compromise-block rate, legitimate-reject rate) operating point, plotted on the same plane as the base study's four published recommendation-time defense points, which are placed as fixed reference markers --- not re-run, since re-running them is not this study's contribution --- and every statistic is reported per cell, never pooled.

\section{Experiments}

\subsection{Experimental Setup}

\textbf{Instruments and source material.} The testbed is a novel multi-step recommend-then-install trajectory harness embedding decisions from the base study's two settings --- standalone-LLM recommendation and agent-harness recommend-then-install \citep{hallucskills} --- over executable multi-step tasks in the style of a real-world executable-API tool-use suite \citep{toolllm}. The registry is a trust-tiered snapshot: a real marketplace's publisher-verification metadata where one is available, otherwise a constructed synthetic registry that exposes the three trust tiers --- deployed analogues of such metadata exist in package ecosystems (npm provenance attestations, PyPI Trusted Publishers), so the three-tier construct models real, not hypothetical, registry practice. The squatting adversary is seeded from the base study's confirmed recurring hallucinated-name set; a held-out probe set of real, useful skills from unverified/independent publishers measures the usability-cost axis. Because the harness, registry snapshots, and probe set are released artifacts, their construction is documented and their contents held disjoint from any tuning split.

\textbf{Replication gate (pre-run pass criterion, not a research question).} Before any gate arm is scored, the harness must reproduce a skill-hallucination rate within the base study's published 36--43\% range on at least one of its own model-and-harness configurations, using its prompt-construction method as closely as feasible. This is the placebo-equivalent for this study: a harness that elicits far less hallucination would leave every squatting arm underpowered and every operating point incomparable to the published points. Failure reopens the harness rather than being reported as a substantive result.

\textbf{Arms.} Gate-off (no intervention) is the control for every gate-on claim. The provenance gate (immediate, $\delta = 0$) is the intervention. A bare-existence gate (allow iff the name resolves in the registry) is the falsification placebo. A delayed-gate arm ($\delta = k$) isolates timing. The base study's four published recommendation-time defense operating points are placed on the plane as fixed reference markers, not re-run.

\textbf{Metrics.} Compromise-block rate (fraction of squatted-install attempts on recurrence-seeded names blocked, per cell); legitimate-unverified-publisher reject rate, reported as the bounded interval [block-only, block-plus-review]; detection lag (step index of first catch; ``uncaught'' when a squatted skill executes and the task still succeeds); recovery cost (extra steps and tokens between an untrusted install and task resolution); reproduced baseline hallucination rate (the replication gate).

\textbf{Sample size and correction.} Each cell is sized at roughly 300 trajectory episodes per arm per model, powered to detect a 10-percentage-point difference in compromise-block rate at 80\% power, two-sided $\alpha = 0.05$ under a two-proportion test; the base rate is re-estimated from the pilot before scale-up. Primary contrasts use Fisher's exact test per cell with Holm--Bonferroni correction across cells. RQ4 is the one contrast whose decision boundary (2 percentage points) lies below this design's 10-point minimum detectable effect: it is therefore pre-registered as an equivalence claim, tested with two one-sided tests (TOST) at a 2-point margin and sized separately --- at a block rate near 90\%, on the order of 3{,}000 episodes per arm (design arithmetic from the stated margin, $\alpha$, and power, to be re-estimated once the pilot pins the base rate); if that cell cannot be run at full size, RQ4 is reported descriptively with its confidence interval, not as a passed equivalence test. Any human-coded disputed rejection is double-annotated with a Cohen's $\kappa \geq 0.7$ acceptance threshold. Every quantity stated in the research questions below is a pre-registered design target, not a measured result.

\textbf{Research questions:}
\begin{enumerate}[leftmargin=1.6em]
\item RQ1 --- Does the squatting attack path execute end-to-end in the harness under no gate?
\item RQ2 --- Does a bare-existence gate fail to block squatted installs?
\item RQ3 --- Does the provenance gate block squatted installs, using zero model judgment?
\item RQ4 --- Does provenance alone suffice --- does adding a model judgment raise the block rate?
\item RQ5 --- What is the usability cost, and is it materially below the published recall collapse?
\item RQ6 --- Does the security outcome worsen monotonically as interception is delayed?
\item RQ7 --- Is the gate's benefit larger where a real install step exists to intercept?
\end{enumerate}

\subsubsection{RQ1: Does the squatting attack path execute end-to-end in the harness under no gate?}
\textbf{Setup.} Under gate-off, run recurrence-seeded squatting episodes and read from sandbox state whether the squatted skill was installed and executed. This establishes the phenomenon is present in the harness before any defense is measured, extending the base study's single proof-of-concept to a systematic rate.

\textbf{Expected outcome.} The compromise-completion rate under gate-off is materially above zero --- at least 20 percentage points --- on recurrence-seeded squatted names, confirming the attack path executes end-to-end in the harness, not just in the base study's single PoC.

\textbf{Falsification condition.} Refuted if the gate-off compromise-completion rate is below 20\% (the attack does not reliably execute in this harness), which would mean the testbed does not reproduce the demonstrated attack and every downstream block-rate comparison is uninterpretable --- reopen the harness rather than score the gate.

\textbf{Hypothesis (plain language).} The demonstrated squatting attack is reliably present in this multi-step testbed, not a one-off.

\textbf{Metrics.} Compromise-completion rate under gate-off, per cell.

\textbf{Result.} \ph

\subsubsection{RQ2: Does a bare-existence gate fail to block squatted installs?}
\textbf{Setup.} Replace the provenance gate with a bare-existence gate (allow iff the recommended name resolves in the registry) and measure its compromise-block rate on the same recurrence-seeded squatted installs. This is the falsification control for the paper's central claim that provenance, not existence, is what does the work: if a bare-existence check sufficed, the whole motivation would collapse.

\textbf{Expected outcome.} The bare-existence gate blocks under 5\% of squatted-malicious installs, because the attacker's entire move is to make the name resolve --- establishing that existence-checking is not a defense against the base study's own attack.

\textbf{Falsification condition.} Refuted if the existence gate blocks 5\% or more of squatted installs, which would mean name existence carries real safety signal after all and the provenance framing is not necessary --- a result that would redirect the paper toward characterizing when existence suffices.

\textbf{Hypothesis (plain language).} Checking that a skill by this name exists provides essentially no protection against squatting, because squatting makes the name exist.

\textbf{Metrics.} Compromise-block rate of the bare-existence gate, per cell; contrast against the provenance gate's block rate (RQ3).

\textbf{Result.} \ph

\subsubsection{RQ3: Does the provenance gate block squatted installs, using zero model judgment?}
\textbf{Setup.} Under the immediate provenance gate ($\delta = 0$), measure the compromise-block rate on recurrence-seeded squatted installs, per cell, confirming that no model call enters the verdict path. A squatted skill held for human-review counts as blocked (no autonomous auto-install), per the gate definition. This RQ is an integrity check on the instrument rather than a finding about the world --- the gate blocks the unknown tier deterministically, so its informative content is the count of leaks; the study's informative endpoints are RQ5--RQ7.

\textbf{Expected outcome.} The provenance gate blocks at least 90\% of squatted-malicious installs using zero model calls, per model-and-harness cell, because every recurrence-seeded squatted skill occupies the unknown-publisher tier by construction and the gate blocks that tier deterministically.

\textbf{Falsification condition.} Refuted if the block rate falls below 90\% in a majority of cells, which --- since squatted skills are unknown-tier by construction --- would indicate a leak in the gate or the snapshot (e.g., a squatted publisher wrongly resolving to a known tier) rather than a property of provenance gating; diagnose and reopen rather than report a low block rate as a finding.

\textbf{Hypothesis (plain language).} A gate that blocks unknown-publisher installs stops almost all squatting, without asking any model whether a name looks real.

\textbf{Metrics.} Compromise-block rate at $\delta = 0$, per cell; count of model calls in the verdict path (target: zero).

\textbf{Result.} \ph

\subsubsection{RQ4: Does provenance alone suffice --- does adding a model judgment raise the block rate?}
\textbf{Setup.} Ablate the provenance-only constraint by adding an LLM ``does this name look real'' judgment to the verdict path (block if either provenance or the model flags the skill) and compare the block rate against the provenance-only gate on the same squatted installs. Because the 2-point boundary is below the study's 10-point powered effect, this contrast runs as the separately sized TOST equivalence cell described in the Experimental Setup.

\textbf{Expected outcome.} Adding the model judgment raises the block rate by less than 2 percentage points over the provenance-only gate --- no material gain --- consistent with the base study's finding that self- and cross-model auditing is coin-flip (51--53\%): the model contributes essentially no reliable signal the provenance tier does not already carry.

\textbf{Falsification condition.} Refuted if adding the model judgment raises the block rate by 2 percentage points or more, which would mean the model contributes real signal beyond provenance and the provenance-only design is leaving protection on the table --- a result that would motivate a hybrid gate rather than a provenance-only one.

\textbf{Hypothesis (plain language).} The model's judgment adds nothing the publisher tier does not already provide, so keeping the model out of the gate costs no protection.

\textbf{Metrics.} Block-rate difference (provenance-plus-LLM minus provenance-only), per cell.

\textbf{Result.} \ph

\subsubsection{RQ5: What is the usability cost, and is it materially below the published recall collapse?}
\textbf{Setup.} Run the held-out probe set of legitimate skills from unverified publishers through the immediate provenance gate and measure the rejection rate, reported as the bounded interval [block-only, block-plus-review]. Compare directly against the base study's published recall-loss figures for its recommendation-time defenses (roughly 84\% real-skill rejection under the strongest, 55\% under a weaker one). Disputed rejections are double-annotated ($\kappa \geq 0.7$).

\textbf{Expected outcome.} The gate's legitimate-unverified-publisher rejection rate is materially below the base study's worst recommendation-time recall loss ($\approx$84\%), specifically below half of it --- under 42 percentage points at the conservative (block-plus-review) end --- because moving the intervention off the recommendation distribution does not penalize a correctly-named real skill for being training-atypical. The rejection rate is not expected to be zero, since a legitimate first-time publisher is indistinguishable from a squatter on provenance alone until trust accrues.

\textbf{Falsification condition.} Refuted if the conservative-end rejection rate on legitimate skills is 42\% or higher --- statistically indistinguishable from the base study's own recommendation-time recall loss --- which would mean moving the intervention point bought no usability improvement and the tradeoff was relocated, not solved.

\textbf{Hypothesis (plain language).} The gate wrongly rejects far fewer real skills than the recommendation-time defenses do, though it cannot spare every unknown first-time publisher.

\textbf{Metrics.} Rejection rate on the legitimate probe set, as the interval [block-only, block-plus-review], per cell; comparison to the published recall-loss figures.

\textbf{Result.} \ph

\subsubsection{RQ6: Does the security outcome worsen monotonically as interception is delayed?}
\textbf{Setup.} Across the delayed-gate arm's timing levels \{immediate ($\delta = 0$), delayed-by-$k$, never (gate-off)\}, holding the trajectory and task difficulty fixed, measure the compromise-execution rate and recovery cost, and test for a monotone trend.

\textbf{Expected outcome.} Compromise-execution rate and recovery cost worsen monotonically as the check is delayed across \{immediate, delayed-$k$, never\} (Jonckheere--Terpstra trend test, $p < 0.05$), holding task difficulty fixed --- later interception means more squatted skills execute and each costs more steps to recover from.

\textbf{Falsification condition.} Refuted if the trend is non-monotone or non-significant ($p \geq 0.05$), which would mean interception timing does not causally govern the security outcome and the ``when it is caught matters'' motivation is not supported --- the finding would then be that only whether, not when, the gate fires matters.

\textbf{Hypothesis (plain language).} The later an untrusted install is caught, the more damage executes and the costlier the recovery --- timing is causal, not incidental.

\textbf{Metrics.} Compromise-execution rate and recovery cost per timing level; Jonckheere--Terpstra trend statistic.

\textbf{Result.} \ph

\subsubsection{RQ7: Is the gate's benefit larger where a real install step exists to intercept?}
\textbf{Setup.} Compare the gate's benefit --- the compromise-block-rate gain of gate-on over gate-off --- between the agent-harness recommend-then-install setting and the standalone-LLM recommendation setting, matching the base study's own standalone-versus-agent split for direct comparability. To keep the contrast well-defined, the standalone arm's recommendation is consumed by a scripted, out-of-loop installer (mirroring the base study's proof-of-concept path from recommendation to install), so the gate can veto in both settings; the contrast then measures the value of in-the-loop interception with recovery versus a post-hoc veto with none.

\textbf{Expected outcome.} The gate's benefit is larger in the agent-harness recommend-then-install setting than in the standalone-LLM recommendation setting, by at least 10 percentage points, because only the former embeds the install inside a live trajectory where interception preserves the task and in-context recovery is possible; the standalone arm's scripted install can be vetoed but not recovered from in-context.

\textbf{Falsification condition.} Refuted if the gate's benefit is statistically indistinguishable between the two settings, which would mean the install step is not where the gate's value comes from --- a result that would broaden the claim from ``install-time interception'' to something the recommendation setting shares.

\textbf{Hypothesis (plain language).} The gate helps most exactly where there is a real auto-install to intercept.

\textbf{Metrics.} Compromise-block-rate gain (gate-on minus gate-off) per setting; between-setting difference.

\textbf{Result.} \ph

\vspace{0.5em}
\noindent\textbf{Further evaluation.} With more time, a natural extension is a single simulated adversarial-retargeting sweep --- does a deployed provenance gate push attackers away from squatting raw hallucinations and toward typosquatting the names of verified publishers instead? This is explicitly a scale-up stretch goal, not a headline result: it is a longer-horizon question that a single simulated sweep can only partially answer, and a live, iteratively-adapting red-team adversary is out of scope for this design.

\subsection{Threats to Validity}

\textbf{Construct validity.} The compromise and usability constructs are only as sound as the trust-tier snapshot and the gate's verdict definition. Ground truth for ``compromise'' is the pre-committed trust-tier snapshot plus sandbox execution state, never a model judgment, so the security label does not inherit the coin-flip detectability the base study documents. The one construct decision that materially moves a headline number is the treatment of the gate's human-review verdict on the reject-rate axis; this is why the usability cost is reported as a bounded interval [block-only, block-plus-review] rather than a single number that would hide the human-review policy, with the conservative end counting every human-review of a legitimate skill as a rejection. A separate construct threat is the fidelity of the squatting adversary: seeding it from arbitrary rather than recurrence-confirmed names would overstate the gate's value, which is why the adversary draws only from the base study's confirmed recurring set.

\textbf{Internal validity.} Two confounds threaten a causal reading. First, whether a squatted install executes must be attributable to the gate's timing, not to task difficulty; the delayed-gate control holds the trajectory and task difficulty fixed and varies only $\delta$, and any episode failing static validation (a task with no well-defined completion state) is excluded before any API spend, with the exclusion rate reported rather than silently dropped. Second, the provenance-only claim is undermined if the gate quietly reads the recommended name; the harness asserts that the gate's verdict path contains zero model calls and no name-string lookup (RQ3/RQ4 measure this directly), and the recurring-name set is architecturally confined to the offline adversary and labeler.

\textbf{External validity.} Findings are scoped to sandboxable, two-party recommend-then-install trajectories with a computable trust-tier snapshot, across the model-and-harness cells actually tested (the base study's standalone and agent-harness settings). They do not cover live uncontrolled registries whose publisher-verification metadata may be forgeable or absent, frontier proprietary models beyond those tested, or a live iteratively-adapting adversary --- the single simulated retargeting sweep is a stretch goal, not evidence about a real adaptive attacker. The released harness, registry snapshots, and probe set are the intended vehicle for extension to configurations beyond the tested grid.

\textbf{Conclusion validity.} A single cell's block rate cannot distinguish a real gate effect from harness noise, so every operating-point and timing statistic is reported per model-and-harness cell with Holm--Bonferroni correction across cells, never pooled, and every squatting arm is scored only after the replication gate confirms the harness reproduces the base study's 36--43\% hallucination rate --- the placebo-equivalent whose failure reopens the harness rather than being read as a null result. The sample size is pre-registered to detect a 10-percentage-point block-rate difference at 80\% power, and underpowered cells are pre-declared inconclusive rather than reported as null.

\section{Related Work}

The closest prior work to this paper is the large-scale hallucinated-skill-recommendation study of \citet{hallucskills}, and it is worth stating the delta plainly before surveying the wider landscape. That study establishes, at scale, everything this paper takes as given: that agents hallucinate skill names at 36--43\% across every configuration, that the names recur predictably enough to be squatted, that self- and cross-model auditing detects them no better than a coin flip, that all four of its recommendation-time defenses trade a lower hallucination rate for a collapse in real-skill recall, and --- in a single sandboxed proof of concept --- that the full squatting attack path executes end to end. We concede all of it and build on it. What that study does not do, and explicitly defers, is exactly this paper's contribution: it runs no multi-step trajectory study with a downstream task outcome, measures no detection timing and no recovery cost, evaluates no install-time intervention, and states in its own closing sentence that ``registry-level name reservations and verified recommendation pipelines'' are the needed fix without building either. We move the intervention point from recommendation to install, swap the signal from name existence (which the study's own squatting threat model defeats) to publisher provenance, and measure the operating point, timing curve, and usability cost it never ran.

\textbf{Supply-chain squatting from hallucinated names, and registry-trust proposals.} The attack-precondition pattern predates the skill setting: code-generating LLMs hallucinate package names at scale, creating the ``slopsquatting'' surface in which attackers pre-register hallucinated dependencies \citep{spracklen,phantoms}; a 2026 frontier-cohort replication finds the rate compressed but the threat intact, with a set of model-agnostic hallucinated names still registrable after registries' deployed defenses \citep{packrange}; and prompting-side attacks can stealthily steer agents toward more speculative names \citep{npa}. Defenses in that lineage sit, like the base study's, on the generation side --- decoding-time monitoring against authoritative package lists \citep{packmonitor} and model editing of the package-validity boundary \citep{bound} --- the very intervention point whose structural cost this paper measures an alternative to. On the mechanism side, design-level proposals for tool-registry trust exist: a zero-trust tool registry with admin-controlled registration and dynamic trust scoring \citep{ztr}, and OAuth-based cryptographic tool-identity verification for the Model Context Protocol \citep{etdi}. We concede the mechanism class to these proposals; both are security architectures, and neither reports a measured security/usability operating point against a hallucination-seeded squatting adversary, a usability cost on legitimate unverified publishers, or any interception-timing measurement --- the quantities this paper's testbed is built to produce.

\textbf{Tool-use grounding and verification.} A line of work grounds proposed agent actions in real, executable resources to combat brittleness: a large-scale executable-API tool-use suite trains and evaluates agents against thousands of real APIs with multi-step solution paths \citep{toolllm}, and a tool-augmented reward model invokes real external tools mid-evaluation rather than trusting an internally generated score \citep{themis}. Both treat grounding as the fix for unreliable \emph{reasoning}, but both assume the tool set is fixed, known, and trustworthy by construction --- the exact assumption the base study shows is false for open skill registries --- and neither asks whether a \emph{proposed} resource's publisher can be trusted. The closest mechanism-level precedents for on-demand external checking, Self-RAG's learned decision to retrieve and critique support for a generated claim \citep{selfrag} and Function Masking's training of a policy to abstain when no real function applies within a known set \citep{funcmasking}, both operate by retraining the policy over a known, fixed resource set and never face the existence-or-provenance question a drop-in install-time gate must answer; they are method-side references, not runnable comparators for this study.

\textbf{Failure attribution and recover-early architectures in agent trajectories.} A second line studies how an early error propagates through a multi-step trajectory. Counterfactual-replay failure attribution localizes the earliest decisive error step by holding a trajectory fixed and substituting a corrected action \citep{agentracer}; tree-structured planning with per-node verification and backtracking catches a failing step and recovers before continuing down a linear pipeline \citep{genartist}. These are the closest precedents for ``catch the failure early and recover,'' and this paper's controlled-delay design borrows the counterfactual logic of holding the trajectory fixed while moving the intervention point --- but the prior work diagnoses a \emph{completed} trajectory post hoc, or verifies image-generation quality, and neither treats tool provenance as the thing checked nor measures recovery cost as a function of how early the check fires, which are exactly this paper's timing measurements.

\textbf{Trajectory-centered agent testbeds and data verification.} A third line builds and verifies agent trajectories as instruments. Self-evolving agent benchmarks treat a reproducible test-time trajectory as evidence a task is solvable and convert it into a verified benchmark item \citep{selfevolve}, and ToolACE-MT synthesizes multi-turn tool-use dialogues behind an explicit offline verification stage \citep{toolacemt}. Both apply verification to \emph{construct} or \emph{validate trajectory data}, once and offline; neither runs a live, per-step safety gate on a deployed agent's real-time install proposals, which is the online intervention this paper measures. None of these three lines establishes whether an install-time, provenance-keyed gate can block squatted installs without the usability collapse that every recommendation-time defense in the base study suffers --- the question this paper's testbed is built to answer.

\section{Conclusion}

This work reframes a governance assumption --- that checking a recommended skill against the registry disciplines hallucinated and malicious tool installs --- as an empirical question with a precise structure, and corrects it at two points. The checkpoint belongs at install time, not recommendation time, because a provenance check there need not constrain the recommendation distribution and so need not pay the recall collapse that every recommendation-time defense in the base study pays; and the signal must be publisher provenance, not name existence, because the squatting attack makes the malicious name exist by construction. The central methodological wager is that whether this relocation solves or merely relocates the security/usability tradeoff is decided by one measured number --- the rejection rate on real skills from unverified publishers --- and that whether interception timing matters is decided by a controlled delay that holds the task fixed; both are measured, not assumed, and a bare-existence placebo gate is carried specifically to falsify the claim that provenance rather than existence does the work. Two limitations bound the claim from the outset and are natural next steps rather than afterthoughts: the gate cannot drive the legitimate-first-time-publisher rejection rate to zero on provenance alone until trust accrues, so a trust-accrual mechanism for honest newcomers is the obvious follow-up; and whether a deployed gate simply pushes attackers toward typosquatting verified publishers is a longer-horizon adversarial question this design can only begin to probe with a single simulated sweep. Whichever way the operating point lands relative to the published recommendation-time frontier, the measurement is decision-relevant: it tells a marketplace or harness builder whether install-time provenance is a cheap first line of defense or whether it, too, buys security by rejecting the honest newcomers a healthy ecosystem depends on.

\bibliographystyle{plainnat}
\bibliography{refs}

\vspace{1em}
\hrule
\vspace{1em}

\section*{Pre-handoff checklist}
\begin{itemize}[leftmargin=1.4em]
\item \textbf{Organization/navigation} --- narrative flows problem $\to$ method $\to$ experiments $\to$ related work without forward references to undefined terms; the two corrections (install-time vs recommendation-time; provenance vs existence) are introduced in the Introduction and reused consistently.
\item \textbf{Terminology discipline} --- ``publisher-trust tier,'' ``provenance gate,'' ``squatting,'' ``compromise-block rate,'' ``detection lag,'' ``operating point'' each defined at first use in the Introduction/Problem Definition, used consistently thereafter; no internal stage IDs (\texttt{sp\_*}, \texttt{m\_*}, \texttt{rq\_*}, \texttt{C*}, \texttt{M*}) leaked into prose --- checked directly against the discussion trace.
\item \textbf{Notation with worked example} --- the formal setup defines $(\tau, x_r, n, b, R, p_k, \theta_k, g_\delta, \delta, \mathcal{H}, \mathcal{L})$ and the Problem Definition/Methodology reuse it consistently; no separate worked numeric example is included since no numbers exist yet.
\item \textbf{RQ framing matches outcomes} --- all 7 RQs are graded/quantitative, each answered in its own subsection with an explicit Falsification condition; the study\_type is empirical, so no RQ is framed as ``our gate beats defense X.''
\item \textbf{No typesetting artifacts} --- only explicit placeholder markers in Results, per the drafting convention this is a pre-registered spec, not a report.
\item \textbf{Artifact plan} --- the trajectory harness, trust-tiered registry snapshots, recurrence-seeded adversary, and legitimate-publisher probe set are the natural release artifacts; not yet built, so not claimed as available.
\end{itemize}

\section*{Hand-off}
This is a draft, not a finished paper --- experiments have not run. Natural next steps: prose polish and venue-specific LaTeX formatting once a target venue is chosen, or revising an earlier stage if piloting surfaces a problem (e.g., if the replication gate cannot reproduce the base study's hallucination rate in this harness, that is a walking-skeleton failure that reopens the harness before any gate arm is scored).

\end{document}
